\documentclass[11pt,a4paper]{article}
\usepackage{fontspec,amsmath,amssymb,amsthm,microtype}
\setmainfont{lmroman10-regular.otf}[BoldFont=lmroman10-bold.otf,ItalicFont=lmroman10-italic.otf,BoldItalicFont=lmroman10-bolditalic.otf]
\usepackage{xcolor}
\usepackage[most]{tcolorbox}
\usepackage{mathrsfs}
\providecommand{\lt}{<}
\providecommand{\gt}{>}
\setmonofont{DejaVuSansMono.ttf}[BoldFont=DejaVuSansMono-Bold.ttf,ItalicFont=DejaVuSansMono-Oblique.ttf,BoldItalicFont=DejaVuSansMono-BoldOblique.ttf,Scale=0.86]
\newfontfamily\leanfallback{STIXTwoMath-Regular.otf}[Scale=0.86]
\definecolor{leanrule}{gray}{0.62}
\definecolor{leanback}{gray}{0.962}
\definecolor{leantitle}{gray}{0.22}
\definecolor{leanmuted}{gray}{0.38}
\newtcolorbox{leanbox}[1]{enhanced,breakable,colback=leanback,colframe=leanrule,
  boxrule=0.4pt,arc=0.6mm,left=2.4mm,right=2.4mm,top=1.8mm,bottom=1.8mm,
  before skip=0.8em,after skip=1em,fonttitle=\small\bfseries,coltitle=leantitle,
  colbacktitle=leanback,titlerule=0.3pt,title={#1}}
\newlength{\leancharwidth}
\newenvironment{leancode}{\par\smallskip\begingroup\ttfamily\footnotesize\raggedright
  \settowidth{\leancharwidth}{M}\setlength{\parskip}{0pt}\setlength{\parindent}{0pt}
  \frenchspacing}{\par\endgroup\smallskip}
\newcommand{\leanline}[2]{\hangindent=\dimexpr#1\leancharwidth+4\leancharwidth\relax
  \hangafter=1\noindent\hspace*{#1\leancharwidth}#2\par}
\newcommand{\leansource}[1]{\par{\footnotesize\color{leanmuted}#1}\par}
\newcommand{\leanlabel}[1]{\par\smallskip\noindent{\small\bfseries #1}\par}
\usepackage[margin=25mm]{geometry}
\usepackage[unicode,hidelinks]{hyperref}
\hypersetup{pdfauthor={LangevinLean Project}}
\setlength{\emergencystretch}{3em}
\setlength{\parskip}{0.4em}
\setlength{\parindent}{0pt}
\allowdisplaybreaks[2]
\raggedbottom
\title{\texorpdfstring{Coordinate Gaussian moments and \(L^2\) Brownian continuity}{Coordinate Gaussian moments and L² Brownian continuity}}
\author{LangevinLean Project}
\date{}
\begin{document}\maketitle
Draft mathematical exposition; these notes may contain mistakes.


For a finite coordinate set \(\iota\), scalar laws \(X_i\sim\mathcal N(0,v)\) give \(\mathbb E\|X\|^2=|\iota|v\) by summing squares. The corresponding class has norm \(\sqrt{|\iota|v}\); representative changes on null sets do not alter this calculation. Spatial coordinate independence is not needed for these moment identities.



For Brownian coordinate processes, differences at times \(s,t\) have coordinate variance \(|t-s|\). Therefore \(\|[B_t]-[B_s]\|_2=\sqrt{|\iota|\,|t-s|}\), which proves continuity in mean square. Multiplication by \(\sqrt2\) gives the Langevin noise norm. Extension to all real times uses \(t\mapsto\max(t,0)\) and retains the same positive-time increment formulas.


\section{\texorpdfstring{Square-integrability from Gaussian coordinate laws}{Square-integrability from Gaussian coordinate laws}}


Each centered Gaussian coordinate has a finite second moment. In finite Euclidean dimension, \(\|X\|^2=\sum_iX_i^2\), so the finite sum is integrable. Coordinatewise measurability also assembles to vector measurability, giving \(X\in L^2\). Coordinate independence is unnecessary for this moment statement.



The finite-dimensional norm square is a finite sum of coordinate squares. Each coordinate has a finite Gaussian second moment, proving vector square-integrability without requiring independence for this particular conclusion.

\begin{leanbox}{Lean: \texttt{gaussian\_\allowbreak{}coordinates\_\allowbreak{}memLp}}
\leansource{\texttt{gaussian\_\allowbreak{}coordinates\_\allowbreak{}memLp} \quad UnadjustedLangevin/\allowbreak{}GaussianCoordinates.lean, lines 27–32}
\begin{leancode}
\leanline{0}{theorem\ gaussian\_coordinates\_memLp\ (X\ :\ Ω\ →\ EuclideanSpace\ ℝ\ ι)\ (v\ :\ ℝ≥0)}
\leanline{4}{(hlaw\ :\ ∀\ i\ :\ ι,\ HasLaw\ (fun\ ω\ ↦\ X\ ω\ i)\ (gaussianReal\ 0\ v)\ P)\ :}
\leanline{4}{MemLp\ X\ 2\ P}
\end{leancode}
\leanlabel{How the Lean proof follows the argument}
\texttt{gaussian\_\allowbreak{}coordinates\_\allowbreak{}mem\allowbreak{}Lp}: the source steps apply or rewrite \texttt{mem\allowbreak{}Lp\_\allowbreak{}pi\allowbreak{}Lp\_\allowbreak{}iff.\allowbreak{}mpr}, \texttt{mem\allowbreak{}Lp\_\allowbreak{}id\_\allowbreak{}gaussian\allowbreak{}Real}.
\end{leanbox}

\section{\texorpdfstring{The Euclidean Gaussian second moment}{The Euclidean Gaussian second moment}}


Use \(\|X\|^2=\sum_iX_i^2\) and finite-sum integral linearity. The coordinate law \(X_i\sim\mathcal N(0,v)\) gives \(\mathbb EX_i^2=v\), hence \(\mathbb E\|X\|^2=|\iota|v\). The dimension factor is exactly the cardinality of the finite coordinate set.



Summing the coordinate second moments gives \(dv\). No mixed coordinate products occur, so the same formula remains valid for correlated coordinates having these marginal variances.

\begin{leanbox}{Lean: \texttt{gaussian\_\allowbreak{}coordinates\_\allowbreak{}second\_\allowbreak{}moment}}
\leansource{\texttt{gaussian\_\allowbreak{}coordinates\_\allowbreak{}second\_\allowbreak{}moment} \quad UnadjustedLangevin/\allowbreak{}GaussianCoordinates.lean, lines 35–45}
\begin{leancode}
\leanline{0}{theorem\ gaussian\_coordinates\_second\_moment\ (X\ :\ Ω\ →\ EuclideanSpace\ ℝ\ ι)\ (v\ :\ ℝ≥0)}
\leanline{4}{(hlaw\ :\ ∀\ i\ :\ ι,\ HasLaw\ (fun\ ω\ ↦\ X\ ω\ i)\ (gaussianReal\ 0\ v)\ P)\ :}
\leanline{4}{(∫\ ω,\ ‖X\ ω‖\ \textasciicircum{}\ 2\ ∂P)\ =\ (Fintype.card\ ι\ :\ ℝ)\ *\ (v\ :\ ℝ)}
\end{leancode}
\leanlabel{How the Lean proof follows the argument}
\texttt{gaussian\_\allowbreak{}coordinates\_\allowbreak{}second\_\allowbreak{}moment}: the intermediate facts include \texttt{hmem}, \texttt{hmom}, in source order; the source steps apply or rewrite \texttt{integral\_\allowbreak{}comp}, \texttt{fun\_\allowbreak{}prop}, \texttt{Finset.\allowbreak{}sum\_\allowbreak{}const}, \texttt{Finset.\allowbreak{}card\_\allowbreak{}univ}, \texttt{nsmul\_\allowbreak{}eq\_\allowbreak{}mul}.
\end{leanbox}

\section{\texorpdfstring{The squared norm of the Gaussian class}{The squared norm of the Gaussian class}}


For the given \(L^2\) class, its squared norm is the integral of the squared norm of its representative. Apply the preceding Gaussian coordinate moment formula to obtain \(\|X\|_2^2=|\iota|v\).



The class representative has the same squared-norm integral as the original vector. Thus the \(L^2\) norm square is exactly the already computed total coordinate variance.

\begin{leanbox}{Lean: \texttt{gaussian\_\allowbreak{}coordinates\_\allowbreak{}Lp\_\allowbreak{}norm\_\allowbreak{}sq}}
\leansource{\texttt{gaussian\_\allowbreak{}coordinates\_\allowbreak{}Lp\_\allowbreak{}norm\_\allowbreak{}sq} \quad UnadjustedLangevin/\allowbreak{}GaussianCoordinates.lean, lines 48–53}
\begin{leancode}
\leanline{0}{theorem\ gaussian\_coordinates\_Lp\_norm\_sq\ (X\ :\ Lp\ (EuclideanSpace\ ℝ\ ι)\ 2\ P)\ (v\ :\ ℝ≥0)}
\leanline{4}{(hlaw\ :\ ∀\ i\ :\ ι,\ HasLaw\ (fun\ ω\ ↦\ X\ ω\ i)\ (gaussianReal\ 0\ v)\ P)\ :}
\leanline{4}{‖X‖\ \textasciicircum{}\ 2\ =\ (Fintype.card\ ι\ :\ ℝ)\ *\ (v\ :\ ℝ)}
\end{leancode}
\leanlabel{How the Lean proof follows the argument}
\texttt{gaussian\_\allowbreak{}coordinates\_\allowbreak{}Lp\_\allowbreak{}norm\_\allowbreak{}sq}: the source steps apply or rewrite \texttt{real\_\allowbreak{}inner\_\allowbreak{}self\_\allowbreak{}eq\_\allowbreak{}norm\_\allowbreak{}sq}, \texttt{L2.\allowbreak{}inner\_\allowbreak{}def}.
\end{leanbox}

\section{\texorpdfstring{The norm of the Gaussian class}{The norm of the Gaussian class}}


Take nonnegative square roots in \(\|X\|_2^2=|\iota|v\). Norm nonnegativity selects the nonnegative root, giving \(\|X\|_2=\sqrt{|\iota|v}\), including zero variance and zero dimension.



Taking nonnegative square roots turns the second-moment identity into the root mean square norm. The formula covers zero variance and the empty coordinate set.

\begin{leanbox}{Lean: \texttt{gaussian\_\allowbreak{}coordinates\_\allowbreak{}Lp\_\allowbreak{}norm}}
\leansource{\texttt{gaussian\_\allowbreak{}coordinates\_\allowbreak{}Lp\_\allowbreak{}norm} \quad UnadjustedLangevin/\allowbreak{}GaussianCoordinates.lean, lines 56–60}
\begin{leancode}
\leanline{0}{theorem\ gaussian\_coordinates\_Lp\_norm\ (X\ :\ Lp\ (EuclideanSpace\ ℝ\ ι)\ 2\ P)\ (v\ :\ ℝ≥0)}
\leanline{4}{(hlaw\ :\ ∀\ i\ :\ ι,\ HasLaw\ (fun\ ω\ ↦\ X\ ω\ i)\ (gaussianReal\ 0\ v)\ P)\ :}
\leanline{4}{‖X‖\ =\ Real.sqrt\ ((Fintype.card\ ι\ :\ ℝ)\ *\ (v\ :\ ℝ))}
\end{leancode}
\leanlabel{How the Lean proof follows the argument}
\texttt{gaussian\_\allowbreak{}coordinates\_\allowbreak{}Lp\_\allowbreak{}norm}: the source steps apply or rewrite \texttt{gaussian\_\allowbreak{}coordinates\_\allowbreak{}Lp\_\allowbreak{}norm\_\allowbreak{}sq}.
\end{leanbox}

\section{\texorpdfstring{Coordinate laws after forming the class}{Coordinate laws after forming the class}}


A representative of the square-integrable class \([X]\) equals \(X\) almost everywhere. Coordinate evaluation preserves that equality. Since pushforward laws do not change under almost-everywhere equality, the coordinate Gaussian laws transfer to the class representative.



The coordinate of the selected class representative agrees almost everywhere with the original coordinate variable. Their pushforward laws therefore coincide with the stated scalar Gaussian law.

\begin{leanbox}{Lean: \texttt{gaussian\_\allowbreak{}coordinates\_\allowbreak{}toLp\_\allowbreak{}coordinateLaw}}
\leansource{\texttt{gaussian\_\allowbreak{}coordinates\_\allowbreak{}toLp\_\allowbreak{}coordinateLaw} \quad UnadjustedLangevin/\allowbreak{}GaussianCoordinates.lean, lines 63–69}
\begin{leancode}
\leanline{0}{theorem\ gaussian\_coordinates\_toLp\_coordinateLaw}
\leanline{4}{(X\ :\ Ω\ →\ EuclideanSpace\ ℝ\ ι)\ (v\ :\ ℝ≥0)}
\leanline{4}{(hlaw\ :\ ∀\ i\ :\ ι,\ HasLaw\ (fun\ ω\ ↦\ X\ ω\ i)\ (gaussianReal\ 0\ v)\ P)\ (i\ :\ ι)\ :}
\leanline{4}{HasLaw\ (fun\ ω\ ↦\ (gaussian\_coordinates\_memLp\ P\ X\ v\ hlaw).toLp\ X\ ω\ i)}
\leanline{6}{(gaussianReal\ 0\ v)\ P}
\end{leancode}
\leanlabel{How the Lean proof follows the argument}
\texttt{gaussian\_\allowbreak{}coordinates\_\allowbreak{}to\allowbreak{}Lp\_\allowbreak{}coordinate\allowbreak{}Law}: the concluding step is \texttt{exact (hlaw i).\allowbreak{}congr}.
\end{leanbox}

\section{\texorpdfstring{The norm is independent of the chosen representative}{The norm is independent of the chosen representative}}


The class \([X]\) has the transferred coordinate laws and the same total second moment as \(X\). Hence \(\|[X]\|_2^2=|\iota|v\). Both sides are nonnegative, so taking square roots gives the exact Gaussian class norm.



Changing representatives on a null set leaves the squared norm integral unchanged. The class norm is consequently the same \(\sqrt{dv}\) computed from the original coordinates.

\begin{leanbox}{Lean: \texttt{gaussian\_\allowbreak{}coordinates\_\allowbreak{}toLp\_\allowbreak{}norm}}
\leansource{\texttt{gaussian\_\allowbreak{}coordinates\_\allowbreak{}toLp\_\allowbreak{}norm} \quad UnadjustedLangevin/\allowbreak{}GaussianCoordinates.lean, lines 72–77}
\begin{leancode}
\leanline{0}{theorem\ gaussian\_coordinates\_toLp\_norm\ (X\ :\ Ω\ →\ EuclideanSpace\ ℝ\ ι)\ (v\ :\ ℝ≥0)}
\leanline{4}{(hlaw\ :\ ∀\ i\ :\ ι,\ HasLaw\ (fun\ ω\ ↦\ X\ ω\ i)\ (gaussianReal\ 0\ v)\ P)\ :}
\leanline{4}{‖(gaussian\_coordinates\_memLp\ P\ X\ v\ hlaw).toLp\ X‖\ =}
\leanline{6}{Real.sqrt\ ((Fintype.card\ ι\ :\ ℝ)\ *\ (v\ :\ ℝ))}
\end{leancode}
\leanlabel{How the Lean proof follows the argument}
\texttt{gaussian\_\allowbreak{}coordinates\_\allowbreak{}to\allowbreak{}Lp\_\allowbreak{}norm}: the proof is the application \texttt{gaussian\_\allowbreak{}coordinates\_\allowbreak{}Lp\_\allowbreak{}norm P \_\allowbreak{} v}. It specializes that result with the hypotheses and objects in the exact statement.
\end{leanbox}

\section{\texorpdfstring{The Langevin-scaled Gaussian norm}{The Langevin-scaled Gaussian norm}}


The scalar-norm identity gives \(\|\sqrt2[X]\|_2=\sqrt2\sqrt{|\iota|v}\). Nonnegativity permits the square-root product rules, producing \(\sqrt{2|\iota|}\sqrt v\). This is the precise Langevin noise constant used in the local error estimate.



Norm homogeneity multiplies the Gaussian root mean square norm by the absolute scale factor. For the Langevin scale \(\sqrt2\), the result is \(\sqrt{2d}\sqrt v\).

\begin{leanbox}{Lean: \texttt{gaussian\_\allowbreak{}coordinates\_\allowbreak{}scaled\_\allowbreak{}l2\_\allowbreak{}norm}}
\leansource{\texttt{gaussian\_\allowbreak{}coordinates\_\allowbreak{}scaled\_\allowbreak{}l2\_\allowbreak{}norm} \quad UnadjustedLangevin/\allowbreak{}GaussianCoordinates.lean, lines 80–88}
\begin{leancode}
\leanline{0}{theorem\ gaussian\_coordinates\_scaled\_l2\_norm\ (X\ :\ Ω\ →\ EuclideanSpace\ ℝ\ ι)\ (v\ :\ ℝ≥0)}
\leanline{4}{(hlaw\ :\ ∀\ i\ :\ ι,\ HasLaw\ (fun\ ω\ ↦\ X\ ω\ i)\ (gaussianReal\ 0\ v)\ P)\ :}
\leanline{4}{‖Real.sqrt\ 2\ •\ (gaussian\_coordinates\_memLp\ P\ X\ v\ hlaw).toLp\ X‖\ =}
\leanline{6}{Real.sqrt\ (2\ *\ (Fintype.card\ ι\ :\ ℝ))\ *\ Real.sqrt\ (v\ :\ ℝ)}
\end{leancode}
\leanlabel{How the Lean proof follows the argument}
\texttt{gaussian\_\allowbreak{}coordinates\_\allowbreak{}scaled\_\allowbreak{}l2\_\allowbreak{}norm}: the source steps apply or rewrite \texttt{norm\_\allowbreak{}smul}, \texttt{Real.\allowbreak{}norm\_\allowbreak{}eq\_\allowbreak{}abs}, \texttt{abs\_\allowbreak{}of\_\allowbreak{}nonneg}, \texttt{Real.\allowbreak{}sqrt\_\allowbreak{}nonneg}. The remaining scalar rearrangements are checked by \texttt{ring}.
\end{leanbox}

\section{\texorpdfstring{Assembling Brownian coordinates}{Assembling Brownian coordinates}}


Assemble the scalar coordinate processes into \(B_t(\omega)=(W_i(t,\omega))_{i\in\iota}\). Finite Euclidean coordinates determine the vector; no new random source or independence is introduced by this definition.



The vector at time \(t\) is assembled from the spatial scalar processes at that same time. Coordinate evaluation recovers each scalar process directly.

\begin{leanbox}{Lean: \texttt{brownianCoordinates}}
\leansource{\texttt{brownianCoordinates} \quad UnadjustedLangevin/\allowbreak{}GaussianCoordinates.lean, lines 91–93}
\begin{leancode}
\leanline{0}{def\ brownianCoordinates\ (W\ :\ ι\ →\ ℝ≥0\ →\ Ω\ →\ ℝ)\ (t\ :\ ℝ≥0)\ :}
\leanline{4}{Ω\ →\ EuclideanSpace\ ℝ\ ι\ :=}
\leanline{2}{fun\ ω\ ↦\ WithLp.toLp\ 2\ (fun\ i\ ↦\ W\ i\ t\ ω)}
\end{leancode}
\leanlabel{How the Lean proof follows the argument}
\texttt{brownian\allowbreak{}Coordinates} is a definition with defining expression \texttt{fun \ensuremath{\omega} \ensuremath{\mapsto} With\allowbreak{}Lp.\allowbreak{}to\allowbreak{}Lp 2 (fun i \ensuremath{\mapsto} W i t \ensuremath{\omega})}. The right-hand side specifies the object used in the mathematical statement.
\end{leanbox}

\section{\texorpdfstring{The coordinate marginal laws}{The coordinate marginal laws}}


At a fixed time \(t\ge0\), each scalar pre-Brownian coordinate has law \(\mathcal N(0,t)\). Unfolding the assembled vector shows that its \(i\)-th coordinate is exactly that scalar variable, giving the required coordinate law.



Each coordinate process has Brownian marginal law \(\mathcal N(0,t)\). The vector construction changes only the codomain, so these fixed-time coordinate laws are unchanged.

\begin{leanbox}{Lean: \texttt{brownianCoordinates\_\allowbreak{}coordinateLaw}}
\leansource{\texttt{brownianCoordinates\_\allowbreak{}coordinateLaw} \quad UnadjustedLangevin/\allowbreak{}GaussianCoordinates.lean, lines 97–100}
\begin{leancode}
\leanline{0}{theorem\ brownianCoordinates\_coordinateLaw\ (W\ :\ ι\ →\ ℝ≥0\ →\ Ω\ →\ ℝ)}
\leanline{4}{(hW\ :\ ∀\ i,\ IsPreBrownianReal\ (W\ i)\ P)\ (t\ :\ ℝ≥0)\ (i\ :\ ι)\ :}
\leanline{4}{HasLaw\ (fun\ ω\ ↦\ brownianCoordinates\ W\ t\ ω\ i)\ (gaussianReal\ 0\ t)\ P}
\end{leancode}
\leanlabel{How the Lean proof follows the argument}
\texttt{brownian\allowbreak{}Coordinates\_\allowbreak{}coordinate\allowbreak{}Law}: the proof is the application \texttt{(h\allowbreak{}W i).\allowbreak{}has\allowbreak{}Law\_\allowbreak{}eval t}. It specializes that result with the hypotheses and objects in the exact statement.
\end{leanbox}

\section{\texorpdfstring{A vector Brownian increment}{A vector Brownian increment}}


Define the vector increment \(\Delta_aB_t=B_{a+t}-B_a\) coordinate by coordinate. This records both the starting time and elapsed nonnegative time, matching the local noise input on an arbitrary Langevin segment.



Subtracting the vector at the starting time yields the vector of scalar increments. A segment beginning at the global time uses this increment as its local additive input.

\begin{leanbox}{Lean: \texttt{brownianCoordinatesIncrement}}
\leansource{\texttt{brownianCoordinatesIncrement} \quad UnadjustedLangevin/\allowbreak{}GaussianCoordinates.lean, lines 103–105}
\begin{leancode}
\leanline{0}{def\ brownianCoordinatesIncrement\ (W\ :\ ι\ →\ ℝ≥0\ →\ Ω\ →\ ℝ)\ (a\ t\ :\ ℝ≥0)\ :}
\leanline{4}{Ω\ →\ EuclideanSpace\ ℝ\ ι\ :=}
\leanline{2}{fun\ ω\ ↦\ brownianCoordinates\ W\ (a\ +\ t)\ ω\ -\ brownianCoordinates\ W\ a\ ω}
\end{leancode}
\leanlabel{How the Lean proof follows the argument}
\texttt{brownian\allowbreak{}Coordinates\allowbreak{}Increment} is a definition with defining expression \texttt{fun \ensuremath{\omega} \ensuremath{\mapsto} brownian\allowbreak{}Coordinates W (a + t) \ensuremath{\omega} -\allowbreak{} brownian\allowbreak{}Coordinates W a \ensuremath{\omega}}. The right-hand side specifies the object used in the mathematical statement.
\end{leanbox}

\section{\texorpdfstring{The coordinate increment laws}{The coordinate increment laws}}


The scalar Brownian increment law depends only on elapsed time: \(W_i(a+t)-W_i(a)\sim\mathcal N(0,t)\). Simplify \((a+t)-a=t\) in the variance and unfold the vector increment to obtain its coordinate law.



The scalar increment law has variance equal to elapsed time, independent of starting time. Thus the same coordinate Gaussian moment argument applies to every segment.

\begin{leanbox}{Lean: \texttt{brownianCoordinatesIncrement\_\allowbreak{}coordinateLaw}}
\leansource{\texttt{brownianCoordinatesIncrement\_\allowbreak{}coordinateLaw} \quad UnadjustedLangevin/\allowbreak{}GaussianCoordinates.lean, lines 109–118}
\begin{leancode}
\leanline{0}{theorem\ brownianCoordinatesIncrement\_coordinateLaw\ (W\ :\ ι\ →\ ℝ≥0\ →\ Ω\ →\ ℝ)}
\leanline{4}{(hW\ :\ ∀\ i,\ IsPreBrownianReal\ (W\ i)\ P)\ (a\ t\ :\ ℝ≥0)\ (i\ :\ ι)\ :}
\leanline{4}{HasLaw\ (fun\ ω\ ↦\ brownianCoordinatesIncrement\ W\ a\ t\ ω\ i)\ (gaussianReal\ 0\ t)\ P}
\end{leancode}
\leanlabel{How the Lean proof follows the argument}
\texttt{brownian\allowbreak{}Coordinates\allowbreak{}Increment\_\allowbreak{}coordinate\allowbreak{}Law}: the source steps apply or rewrite \texttt{has\allowbreak{}Law\_\allowbreak{}sub}, \texttt{NNReal.\allowbreak{}eq}, \texttt{add\_\allowbreak{}sub\_\allowbreak{}cancel\_\allowbreak{}left}, \texttt{abs\_\allowbreak{}of\_\allowbreak{}nonneg}, \texttt{t.\allowbreak{}coe\_\allowbreak{}nonneg}.
\end{leanbox}

\section{\texorpdfstring{Square-integrability of the vector increment}{Square-integrability of the vector increment}}


Apply the finite sum of the Gaussian coordinate second moments to the increment's coordinate laws of variance \(t\). The assembled Brownian vector increment is square integrable even without independence between different coordinate processes.



Every increment coordinate has finite Gaussian second moment. The finite sum criterion then establishes that the whole vector increment belongs to \(L^2\).

\begin{leanbox}{Lean: \texttt{brownianCoordinatesIncrement\_\allowbreak{}memLp}}
\leansource{\texttt{brownianCoordinatesIncrement\_\allowbreak{}memLp} \quad UnadjustedLangevin/\allowbreak{}GaussianCoordinates.lean, lines 121–125}
\begin{leancode}
\leanline{0}{theorem\ brownianCoordinatesIncrement\_memLp\ (W\ :\ ι\ →\ ℝ≥0\ →\ Ω\ →\ ℝ)}
\leanline{4}{(hW\ :\ ∀\ i,\ IsPreBrownianReal\ (W\ i)\ P)\ (a\ t\ :\ ℝ≥0)\ :}
\leanline{4}{MemLp\ (brownianCoordinatesIncrement\ W\ a\ t)\ 2\ P}
\end{leancode}
\leanlabel{How the Lean proof follows the argument}
\texttt{brownian\allowbreak{}Coordinates\allowbreak{}Increment\_\allowbreak{}mem\allowbreak{}Lp}: the proof is the application \texttt{gaussian\_\allowbreak{}coordinates\_\allowbreak{}mem\allowbreak{}Lp P \_\allowbreak{} t}. It specializes that result with the hypotheses and objects in the exact statement.
\end{leanbox}

\section{\texorpdfstring{The scaled vector-increment norm}{The scaled vector-increment norm}}


Apply the exact scaled Gaussian-coordinate norm theorem with variance \(t\). It yields \(\|\sqrt2[\Delta_aB_t]\|_2=\sqrt{2|\iota|}\sqrt t\), giving the norm of the noise from the Brownian marginal law.



Its squared Langevin-scaled norm is \(2dt\). Taking a square root gives a uniform elapsed-time modulus, with no dependence on the segment's absolute starting time.

\begin{leanbox}{Lean: \texttt{brownianCoordinatesIncrement\_\allowbreak{}scaled\_\allowbreak{}l2\_\allowbreak{}norm}}
\leansource{\texttt{brownianCoordinatesIncrement\_\allowbreak{}scaled\_\allowbreak{}l2\_\allowbreak{}norm} \quad UnadjustedLangevin/\allowbreak{}GaussianCoordinates.lean, lines 129–135}
\begin{leancode}
\leanline{0}{theorem\ brownianCoordinatesIncrement\_scaled\_l2\_norm\ (W\ :\ ι\ →\ ℝ≥0\ →\ Ω\ →\ ℝ)}
\leanline{4}{(hW\ :\ ∀\ i,\ IsPreBrownianReal\ (W\ i)\ P)\ (a\ t\ :\ ℝ≥0)\ :}
\leanline{4}{‖Real.sqrt\ 2\ •\ (brownianCoordinatesIncrement\_memLp\ P\ W\ hW\ a\ t).toLp}
\leanline{6}{(brownianCoordinatesIncrement\ W\ a\ t)‖\ =}
\leanline{6}{Real.sqrt\ (2\ *\ (Fintype.card\ ι\ :\ ℝ))\ *\ Real.sqrt\ (t\ :\ ℝ)}
\end{leancode}
\leanlabel{How the Lean proof follows the argument}
\texttt{brownian\allowbreak{}Coordinates\allowbreak{}Increment\_\allowbreak{}scaled\_\allowbreak{}l2\_\allowbreak{}norm}: the proof is the application \texttt{gaussian\_\allowbreak{}coordinates\_\allowbreak{}scaled\_\allowbreak{}l2\_\allowbreak{}norm P \_\allowbreak{} t}. It specializes that result with the hypotheses and objects in the exact statement.
\end{leanbox}

\section{\texorpdfstring{Brownian values in the square-integrable space}{Brownian values in the square-integrable space}}


Convert the assembled Brownian vector at time \(t\) to its \(L^2\) class using its Gaussian-coordinate square-integrability proof. This places the time-indexed process in the Hilbert space where the deterministic integral equation is solved.



The vector at each nonnegative time determines a square-integrable class. The resulting curve takes values in the Hilbert space used for the Bochner integral equation.

\begin{leanbox}{Lean: \texttt{brownianCoordinatesLp}}
\leansource{\texttt{brownianCoordinatesLp} \quad UnadjustedLangevin/\allowbreak{}GaussianCoordinates.lean, lines 138–142}
\begin{leancode}
\leanline{0}{def\ brownianCoordinatesLp\ (W\ :\ ι\ →\ ℝ≥0\ →\ Ω\ →\ ℝ)}
\leanline{4}{(hW\ :\ ∀\ i,\ IsPreBrownianReal\ (W\ i)\ P)\ (t\ :\ ℝ≥0)\ :}
\leanline{4}{Lp\ (EuclideanSpace\ ℝ\ ι)\ 2\ P\ :=}
\leanline{2}{(gaussian\_coordinates\_memLp\ P\ (brownianCoordinates\ W\ t)\ t}
\leanline{4}{(brownianCoordinates\_coordinateLaw\ P\ W\ hW\ t)).toLp\ (brownianCoordinates\ W\ t)}
\end{leancode}
\leanlabel{How the Lean proof follows the argument}
\texttt{brownian\allowbreak{}Coordinates\allowbreak{}Lp} is a definition with defining expression \texttt{(gaussian\_\allowbreak{}coordinates\_\allowbreak{}mem\allowbreak{}Lp P (brownian\allowbreak{}Coordinates W t) t}. The right-hand side specifies the object used in the mathematical statement.
\end{leanbox}

\section{\texorpdfstring{The Brownian representative identity}{The Brownian representative identity}}


For each fixed time, a representative of \([B_t]\) equals the original Brownian vector \(B_t\) almost everywhere. This identifies its coordinate laws while allowing continuity to be studied in the square-mean norm. The exceptional set may depend on that fixed time, which does not affect any marginal or moment statement.



The representative identity permits fixed-time laws and coordinate moments to be evaluated using the original assembled Brownian vector. All equality claims are almost everywhere at the specified time.

\begin{leanbox}{Lean: \texttt{brownianCoordinatesLp\_\allowbreak{}coe}}
\leansource{\texttt{brownianCoordinatesLp\_\allowbreak{}coe} \quad UnadjustedLangevin/\allowbreak{}GaussianCoordinates.lean, lines 145–149}
\begin{leancode}
\leanline{0}{theorem\ brownianCoordinatesLp\_coe\ (W\ :\ ι\ →\ ℝ≥0\ →\ Ω\ →\ ℝ)}
\leanline{4}{(hW\ :\ ∀\ i,\ IsPreBrownianReal\ (W\ i)\ P)\ (t\ :\ ℝ≥0)\ :}
\leanline{4}{brownianCoordinatesLp\ P\ W\ hW\ t\ =ᵐ[P]\ brownianCoordinates\ W\ t}
\end{leancode}
\leanlabel{How the Lean proof follows the argument}
\texttt{brownian\allowbreak{}Coordinates\allowbreak{}Lp\_\allowbreak{}coe}: the proof is the application \texttt{(gaussian\_\allowbreak{}coordinates\_\allowbreak{}mem\allowbreak{}Lp P (brownian\allowbreak{}Coordinates W t) t}. It specializes that result with the hypotheses and objects in the exact statement.
\end{leanbox}

\section{\texorpdfstring{The coordinate law of a class difference}{The coordinate law of a class difference}}


The difference of two selected representatives agrees almost everywhere with \(B_t-B_s\). Brownian increment laws, including the sign reversal when \(t<s\), give a centered Gaussian with variance \(|t-s|\); centered Gaussian symmetry handles the reversed order.



The difference of two class values agrees almost everywhere with the vector of scalar Brownian increments. Its coordinate law therefore has variance equal to the distance between the two times.

\begin{leanbox}{Lean: \texttt{brownianCoordinatesLp\_\allowbreak{}sub\_\allowbreak{}coordinateLaw}}
\leansource{\texttt{brownianCoordinatesLp\_\allowbreak{}sub\_\allowbreak{}coordinateLaw} \quad UnadjustedLangevin/\allowbreak{}GaussianCoordinates.lean, lines 153–163}
\begin{leancode}
\leanline{0}{theorem\ brownianCoordinatesLp\_sub\_coordinateLaw\ (W\ :\ ι\ →\ ℝ≥0\ →\ Ω\ →\ ℝ)}
\leanline{4}{(hW\ :\ ∀\ i,\ IsPreBrownianReal\ (W\ i)\ P)\ (t\ s\ :\ ℝ≥0)\ (i\ :\ ι)\ :}
\leanline{4}{HasLaw\ (fun\ ω\ ↦\ (brownianCoordinatesLp\ P\ W\ hW\ t\ -}
\leanline{6}{brownianCoordinatesLp\ P\ W\ hW\ s)\ ω\ i)}
\leanline{6}{(gaussianReal\ 0\ (nndist\ (t\ :\ ℝ)\ (s\ :\ ℝ)))\ P}
\end{leancode}
\leanlabel{How the Lean proof follows the argument}
\texttt{brownian\allowbreak{}Coordinates\allowbreak{}Lp\_\allowbreak{}sub\_\allowbreak{}coordinate\allowbreak{}Law}: the source steps apply or rewrite \texttt{has\allowbreak{}Law\_\allowbreak{}sub}, \texttt{Pi.\allowbreak{}sub\_\allowbreak{}apply}, \texttt{Pi\allowbreak{}Lp.\allowbreak{}sub\_\allowbreak{}apply}.
\end{leanbox}

\section{\texorpdfstring{The exact norm of a Brownian class difference}{The exact norm of a Brownian class difference}}


Apply the Gaussian \(L^2\) norm formula to the difference class, whose coordinates have variance \(|t-s|\). Thus \(\|[B_t]-[B_s]\|_2=\sqrt{|\iota|\,|t-s|}\). Its limit as the time difference tends to zero gives the following continuity statement.



The norm of a class difference is the square root of the sum of coordinate increment variances. This gives \(\sqrt d\sqrt{|t-s|}\) and immediately tends to zero as \(s\to t\).

\begin{leanbox}{Lean: \texttt{brownianCoordinatesLp\_\allowbreak{}norm\_\allowbreak{}sub}}
\leansource{\texttt{brownianCoordinatesLp\_\allowbreak{}norm\_\allowbreak{}sub} \quad UnadjustedLangevin/\allowbreak{}GaussianCoordinates.lean, lines 166–172}
\begin{leancode}
\leanline{0}{theorem\ brownianCoordinatesLp\_norm\_sub\ (W\ :\ ι\ →\ ℝ≥0\ →\ Ω\ →\ ℝ)}
\leanline{4}{(hW\ :\ ∀\ i,\ IsPreBrownianReal\ (W\ i)\ P)\ (t\ s\ :\ ℝ≥0)\ :}
\leanline{4}{‖brownianCoordinatesLp\ P\ W\ hW\ t\ -\ brownianCoordinatesLp\ P\ W\ hW\ s‖\ =}
\leanline{6}{Real.sqrt\ ((Fintype.card\ ι\ :\ ℝ)\ *\ dist\ (t\ :\ ℝ)\ (s\ :\ ℝ))}
\end{leancode}
\leanlabel{How the Lean proof follows the argument}
\texttt{brownian\allowbreak{}Coordinates\allowbreak{}Lp\_\allowbreak{}norm\_\allowbreak{}sub}: the proof is the application \texttt{simpa only [coe\_\allowbreak{}nndist] using gaussian\_\allowbreak{}coordinates\_\allowbreak{}Lp\_\allowbreak{}norm P}. It specializes that result with the hypotheses and objects in the exact statement.
\end{leanbox}

\section{\texorpdfstring{Continuity in the square-mean norm}{Continuity in the square-mean norm}}


As \(s\to t\), \(|t-s|\to0\), hence the exact norm formula gives \(\|[B_t]-[B_s]\|_2\to0\). The construction uses continuity of the square root and the metric characterization of continuity. This needs finite-dimensional Brownian laws, with no sample-path continuity premise.



The increment norm formula gives a quantitative modulus of continuity in \(L^2\). This supplies the continuous additive input required by the Hilbert-space Langevin existence theorem.

\begin{leanbox}{Lean: \texttt{brownianCoordinatesLp\_\allowbreak{}continuous}}
\leansource{\texttt{brownianCoordinatesLp\_\allowbreak{}continuous} \quad UnadjustedLangevin/\allowbreak{}GaussianCoordinates.lean, lines 176–185}
\begin{leancode}
\leanline{0}{theorem\ brownianCoordinatesLp\_continuous\ (W\ :\ ι\ →\ ℝ≥0\ →\ Ω\ →\ ℝ)}
\leanline{4}{(hW\ :\ ∀\ i,\ IsPreBrownianReal\ (W\ i)\ P)\ :}
\leanline{4}{Continuous\ (brownianCoordinatesLp\ P\ W\ hW)}
\end{leancode}
\leanlabel{How the Lean proof follows the argument}
\texttt{brownian\allowbreak{}Coordinates\allowbreak{}Lp\_\allowbreak{}continuous}: the intermediate fact is \texttt{hc}, in source order; the source steps apply or rewrite \texttt{continuous\_\allowbreak{}iff\_\allowbreak{}continuous\allowbreak{}At.\allowbreak{}mpr}, \texttt{tendsto\_\allowbreak{}iff\_\allowbreak{}norm\_\allowbreak{}sub\_\allowbreak{}tendsto\_\allowbreak{}zero.\allowbreak{}mpr}, \texttt{mul\_\allowbreak{}zero}, \texttt{Real.\allowbreak{}sqrt\_\allowbreak{}zero}, \texttt{hc.\allowbreak{}tendsto}.
\end{leanbox}

\section{\texorpdfstring{The Langevin-scaled increment norm}{The Langevin-scaled increment norm}}


Set the two times in the difference-norm identity to \(a+t\) and \(a\). Their distance is \(t\), and scalar multiplication by \(\sqrt2\) gives \(\sqrt{2|\iota|}\sqrt t\). This is the class-valued version of the earlier increment estimate.



Multiplication of the increment by \(\sqrt2\) produces exactly the diffusion scale of the Langevin equation. Its \(L^2\) norm is \(\sqrt{2d}\sqrt t\) on a segment of duration \(t\).

\begin{leanbox}{Lean: \texttt{brownianCoordinatesLp\_\allowbreak{}increment\_\allowbreak{}scaled\_\allowbreak{}norm}}
\leansource{\texttt{brownianCoordinatesLp\_\allowbreak{}increment\_\allowbreak{}scaled\_\allowbreak{}norm} \quad UnadjustedLangevin/\allowbreak{}GaussianCoordinates.lean, lines 188–199}
\begin{leancode}
\leanline{0}{theorem\ brownianCoordinatesLp\_increment\_scaled\_norm\ (W\ :\ ι\ →\ ℝ≥0\ →\ Ω\ →\ ℝ)}
\leanline{4}{(hW\ :\ ∀\ i,\ IsPreBrownianReal\ (W\ i)\ P)\ (a\ t\ :\ ℝ≥0)\ :}
\leanline{4}{‖Real.sqrt\ 2\ •\ (brownianCoordinatesLp\ P\ W\ hW\ (a\ +\ t)\ -}
\leanline{6}{brownianCoordinatesLp\ P\ W\ hW\ a)‖\ =}
\leanline{6}{Real.sqrt\ (2\ *\ (Fintype.card\ ι\ :\ ℝ))\ *\ Real.sqrt\ (t\ :\ ℝ)}
\end{leancode}
\leanlabel{How the Lean proof follows the argument}
\texttt{brownian\allowbreak{}Coordinates\allowbreak{}Lp\_\allowbreak{}increment\_\allowbreak{}scaled\_\allowbreak{}norm}: the source steps apply or rewrite \texttt{norm\_\allowbreak{}smul}, \texttt{Real.\allowbreak{}norm\_\allowbreak{}eq\_\allowbreak{}abs}, \texttt{add\_\allowbreak{}sub\_\allowbreak{}cancel\_\allowbreak{}left}, \texttt{Real.\allowbreak{}sqrt\_\allowbreak{}mul}, \texttt{Fintype.\allowbreak{}card}. The remaining scalar rearrangements are checked by \texttt{ring}.
\end{leanbox}

\section{\texorpdfstring{The Brownian class at zero}{The Brownian class at zero}}


At time zero every coordinate has the degenerate centered Gaussian law of variance zero. Equivalently its \(L^2\) norm is zero by the moment formula. Norm-zero equivalence proves that the whole class is zero, irrespective of the chosen pointwise representative.



At time zero, each scalar Brownian coordinate has the point-mass law at zero. The assembled vector class has zero \(L^2\) norm and is consequently the zero class.

\begin{leanbox}{Lean: \texttt{brownianCoordinatesLp\_\allowbreak{}zero}}
\leansource{\texttt{brownianCoordinatesLp\_\allowbreak{}zero} \quad UnadjustedLangevin/\allowbreak{}GaussianCoordinates.lean, lines 202–208}
\begin{leancode}
\leanline{0}{@[simp]\ theorem\ brownianCoordinatesLp\_zero\ (W\ :\ ι\ →\ ℝ≥0\ →\ Ω\ →\ ℝ)}
\leanline{4}{(hW\ :\ ∀\ i,\ IsPreBrownianReal\ (W\ i)\ P)\ :}
\leanline{4}{brownianCoordinatesLp\ P\ W\ hW\ 0\ =\ 0}
\end{leancode}
\leanlabel{How the Lean proof follows the argument}
\texttt{brownian\allowbreak{}Coordinates\allowbreak{}Lp\_\allowbreak{}zero}: the source steps apply or rewrite \texttt{norm\_\allowbreak{}eq\_\allowbreak{}zero.\allowbreak{}mp}, \texttt{NNReal.\allowbreak{}coe\_\allowbreak{}zero}, \texttt{mul\_\allowbreak{}zero}, \texttt{Real.\allowbreak{}sqrt\_\allowbreak{}zero}.
\end{leanbox}

\section{\texorpdfstring{A real-time version of the Brownian curve}{A real-time version of the Brownian curve}}


Extend the nonnegative-time \(L^2\) process to real times by \(t\mapsto[B_{\max(t,0)}]\), using the stated nonnegative-time restriction. It is constant zero for negative times and agrees with the original class at nonnegative times.



Extending the nonnegative-time curve to a real interval uses the designated nonnegative time associated with each argument. On the nonnegative horizon this agrees with the original Brownian class.

\begin{leanbox}{Lean: \texttt{brownianCoordinatesRealLp}}
\leansource{\texttt{brownianCoordinatesRealLp} \quad UnadjustedLangevin/\allowbreak{}GaussianCoordinates.lean, lines 212–215}
\begin{leancode}
\leanline{0}{def\ brownianCoordinatesRealLp\ (W\ :\ ι\ →\ ℝ≥0\ →\ Ω\ →\ ℝ)}
\leanline{4}{(hW\ :\ ∀\ i,\ IsPreBrownianReal\ (W\ i)\ P)\ (t\ :\ ℝ)\ :}
\leanline{4}{Lp\ (EuclideanSpace\ ℝ\ ι)\ 2\ P\ :=}
\leanline{2}{brownianCoordinatesLp\ P\ W\ hW\ t.toNNReal}
\end{leancode}
\leanlabel{How the Lean proof follows the argument}
\texttt{brownian\allowbreak{}Coordinates\allowbreak{}Real\allowbreak{}Lp} is a definition with defining expression \texttt{brownian\allowbreak{}Coordinates\allowbreak{}Lp P W h\allowbreak{}W t.\allowbreak{}to\allowbreak{}NNReal}. The right-hand side specifies the object used in the mathematical statement.
\end{leanbox}

\section{\texorpdfstring{Continuity of the real-time curve}{Continuity of the real-time curve}}


The clipping map \(t\mapsto\max(t,0)\) into nonnegative reals is continuous. Compose it with the established \(L^2\)-continuous Brownian map to obtain continuity of the real-time extension, including at zero.



Continuity of the time conversion and of the nonnegative-time \(L^2\) curve gives continuity of the extended curve. This is the real-time domain used by the interval integral.

\begin{leanbox}{Lean: \texttt{brownianCoordinatesRealLp\_\allowbreak{}continuous}}
\leansource{\texttt{brownianCoordinatesRealLp\_\allowbreak{}continuous} \quad UnadjustedLangevin/\allowbreak{}GaussianCoordinates.lean, lines 218–221}
\begin{leancode}
\leanline{0}{theorem\ brownianCoordinatesRealLp\_continuous\ (W\ :\ ι\ →\ ℝ≥0\ →\ Ω\ →\ ℝ)}
\leanline{4}{(hW\ :\ ∀\ i,\ IsPreBrownianReal\ (W\ i)\ P)\ :}
\leanline{4}{Continuous\ (brownianCoordinatesRealLp\ P\ W\ hW)}
\end{leancode}
\leanlabel{How the Lean proof follows the argument}
\texttt{brownian\allowbreak{}Coordinates\allowbreak{}Real\allowbreak{}Lp\_\allowbreak{}continuous}: the proof is the application \texttt{(brownian\allowbreak{}Coordinates\allowbreak{}Lp\_\allowbreak{}continuous P W h\allowbreak{}W).\allowbreak{}comp continuous\_\allowbreak{}real\_\allowbreak{}to\allowbreak{}NNReal}. It specializes that result with the hypotheses and objects in the exact statement.
\end{leanbox}

\section{\texorpdfstring{The real-time curve at zero}{The real-time curve at zero}}


At real time zero the clipping map returns nonnegative zero. The previously proved zero-time \(L^2\) identity therefore simplifies the real-time extension to zero.



At the origin, the converted time is zero, so the extended noise class also vanishes. This verifies the initialization condition in the integral equation.

\begin{leanbox}{Lean: \texttt{brownianCoordinatesRealLp\_\allowbreak{}zero}}
\leansource{\texttt{brownianCoordinatesRealLp\_\allowbreak{}zero} \quad UnadjustedLangevin/\allowbreak{}GaussianCoordinates.lean, lines 223–226}
\begin{leancode}
\leanline{0}{@[simp]\ theorem\ brownianCoordinatesRealLp\_zero\ (W\ :\ ι\ →\ ℝ≥0\ →\ Ω\ →\ ℝ)}
\leanline{4}{(hW\ :\ ∀\ i,\ IsPreBrownianReal\ (W\ i)\ P)\ :}
\leanline{4}{brownianCoordinatesRealLp\ P\ W\ hW\ 0\ =\ 0}
\end{leancode}
\leanlabel{How the Lean proof follows the argument}
\texttt{brownian\allowbreak{}Coordinates\allowbreak{}Real\allowbreak{}Lp\_\allowbreak{}zero}: the source steps apply or rewrite \texttt{Real.\allowbreak{}to\allowbreak{}NNReal\_\allowbreak{}zero}, \texttt{brownian\allowbreak{}Coordinates\allowbreak{}Lp\_\allowbreak{}zero}.
\end{leanbox}

\section{\texorpdfstring{Its scaled increment norm}{Its scaled increment norm}}


For \(a,t\ge0\), clipping leaves \(a\) and \(a+t\) unchanged. Rewrite the real-time extension at both endpoints as the nonnegative-time Brownian classes, then apply their scaled increment formula. This supplies \(\sqrt{2d}\sqrt t\) in the real interval-integral formulation.



Within the nonnegative segment, the real-time version has exactly the original Brownian increment. Its Langevin-scaled norm therefore has the same elapsed-time formula and constant.

\begin{leanbox}{Lean: \texttt{brownianCoordinatesRealLp\_\allowbreak{}increment\_\allowbreak{}scaled\_\allowbreak{}norm}}
\leansource{\texttt{brownianCoordinatesRealLp\_\allowbreak{}increment\_\allowbreak{}scaled\_\allowbreak{}norm} \quad UnadjustedLangevin/\allowbreak{}GaussianCoordinates.lean, lines 230–238}
\begin{leancode}
\leanline{0}{theorem\ brownianCoordinatesRealLp\_increment\_scaled\_norm\ (W\ :\ ι\ →\ ℝ≥0\ →\ Ω\ →\ ℝ)}
\leanline{4}{(hW\ :\ ∀\ i,\ IsPreBrownianReal\ (W\ i)\ P)\ (a\ t\ :\ ℝ)}
\leanline{4}{(ha\ :\ 0\ ≤\ a)\ (ht\ :\ 0\ ≤\ t)\ :}
\leanline{4}{‖Real.sqrt\ 2\ •\ (brownianCoordinatesRealLp\ P\ W\ hW\ (a\ +\ t)\ -}
\leanline{6}{brownianCoordinatesRealLp\ P\ W\ hW\ a)‖\ =}
\leanline{6}{Real.sqrt\ (2\ *\ (Fintype.card\ ι\ :\ ℝ))\ *\ Real.sqrt\ t}
\end{leancode}
\leanlabel{How the Lean proof follows the argument}
\texttt{brownian\allowbreak{}Coordinates\allowbreak{}Real\allowbreak{}Lp\_\allowbreak{}increment\_\allowbreak{}scaled\_\allowbreak{}norm}: the proof is the application \texttt{simpa only [brownian\allowbreak{}Coordinates\allowbreak{}Real\allowbreak{}Lp, Real.\allowbreak{}to\allowbreak{}NNReal\_\allowbreak{}add ha ht,}. It specializes that result with the hypotheses and objects in the exact statement.
\end{leanbox}

\end{document}
